\documentclass[10pt, letterpaper]{article}

\usepackage[T1]{fontenc}
\usepackage{lmodern}
\usepackage[margin=0.75in, top=0.70in, bottom=0.70in]{geometry}
\usepackage{amsmath, amssymb}
\usepackage{booktabs}
\usepackage{hyperref}
\usepackage{fancyvrb}
\usepackage{xcolor}
\usepackage{graphicx}
\usepackage{float}

\setlength{\parskip}{4pt plus 1pt minus 1pt}
\setlength{\parindent}{0pt}
\setlength{\textfloatsep}{8pt}
\setlength{\floatsep}{6pt}
\setlength{\intextsep}{6pt}
\setlength{\abovecaptionskip}{3pt}
\setlength{\belowcaptionskip}{2pt}

\definecolor{codebg}{gray}{0.95}
\DefineVerbatimEnvironment{code}{Verbatim}{%
  fontsize=\fontsize{7.5}{9}\selectfont,
  frame=single, framerule=0.3pt, rulecolor=\color{gray!60},
  fillcolor=\color{codebg}, formatcom=\color{black},
  xleftmargin=0.5em, xrightmargin=0.5em, numbers=left, numbersep=5pt}

% compact list spacing
\let\olditemize\itemize
\renewcommand{\itemize}{\olditemize\setlength{\itemsep}{1pt}\setlength{\topsep}{2pt}}
\let\oldenumerate\enumerate
\renewcommand{\enumerate}{\oldenumerate\setlength{\itemsep}{1pt}\setlength{\topsep}{2pt}}

\hypersetup{colorlinks=true, linkcolor=black, urlcolor=blue}

\title{\vspace{-2.2em}%
  \textbf{Predicting Restaurant Success on the Yelp Open Dataset:}\\[2pt]
  \textbf{A Bayesian Hierarchical and Variational Neural Network Approach}%
  \vspace{-0.4em}}
\author{Brian Butler\\[2pt]
  \small Probabilistic Machine Learning --- Final Project, Spring 2026}
\date{April 2026}

\begin{document}
\maketitle
\vspace{-1em}

%=================================================================
\section*{Abstract}
%=================================================================

This project models \textbf{restaurant success} on the Yelp Open Dataset through two complementary probabilistic lenses: a \textbf{Bayesian hierarchical regression} with partial pooling across cities and cuisines, and a \textbf{variational Bayesian neural network (VBNN)} that fuses structured business metadata with 384-dimensional sentence-transformer embeddings of customer reviews. Success is framed as a multi-objective target combining mean star rating (continuous performance) and the \texttt{is\_open} flag (binary longevity), together with a composite index weighting stars, log review count, and survival. After restricting to five diverse U.S.\ metros---\textbf{Philadelphia, Tampa, Indianapolis, Nashville, and Tucson}---and the top ten restaurant cuisines, the analysis retained \textbf{42,187 restaurants} and a stratified sample of 1.25 million reviews. Ridge, Lasso, and Elastic Net baselines establish a reference RMSE of \textbf{0.682} stars on the held-out test set and a survival classification accuracy of \textbf{0.721} (log-loss 0.541). The hierarchical model, fit with 4-chain NUTS in PyMC~5 under a non-centred parameterisation, attained a test RMSE of \textbf{0.648} with all parameters at $\hat{R} < 1.01$ and \textbf{zero post-warmup divergences}; posterior summaries show that \textbf{price range, outdoor seating, weekend evening hours, and review-count percentile} dominate the fixed effects, while city-level variance ($\tau_{\text{city}} \approx 0.11$) is modestly smaller than cuisine-level variance ($\tau_{\text{cuisine}} \approx 0.17$). The VBNN with mean-field Gaussian variational posterior and Flipout gradient estimator reduced test RMSE to \textbf{0.611} and expected calibration error on survival to \textbf{0.034}, at a $14\times$ compute cost relative to Ridge. Predictive-uncertainty decomposition confirms that epistemic variance concentrates on low-review or newly-opened restaurants, exactly where a point-estimate model would be most misleading. The scientific takeaway is that restaurant success is \textbf{only weakly separable by city or cuisine once operational features are controlled for}, and the PML takeaway is that hierarchical structure earns its complexity on inferential questions while a variational neural net earns its complexity on calibrated prediction---the two models answer different questions and neither dominates the other.

%=================================================================
\section{Overview}
%=================================================================

The Yelp Open Dataset is a large, academically licensed slice of real consumer-review data that Yelp periodically refreshes and distributes for non-commercial research. In the current release it contains roughly \textbf{150,346 businesses, 6,990,280 reviews, 908,915 tips, 1.99 million users, and 131,930 check-in aggregates} spanning eleven metropolitan areas across the United States and Canada, bundled as five line-delimited JSON files totalling 8.65\,GB uncompressed. Businesses carry rich metadata: latitude and longitude, a 47-field attribute object (price tier, parking, Wi-Fi, ambience, outdoor seating, and so on), opening hours, free-text category tags, a half-star average rating, a review count, and a binary \texttt{is\_open} indicator. Reviews carry integer stars, full free text, and engagement counters. This combination of structured attributes and unstructured language makes the dataset unusually well-suited to hybrid probabilistic modelling.

The scientific motivation is straightforward but economically important: \textbf{what makes one restaurant succeed where another fails?} Independent operators face roughly a 17\% first-year closure rate and a 50\% five-year closure rate (Luo and Stark, 2014), and yet the practical intuition about which operational levers matter---price point, menu breadth, hours, neighbourhood, perceived authenticity---rests largely on trade-press anecdote rather than measured effect sizes. Four concrete questions guide the analysis:
\begin{enumerate}
  \item \textbf{Which features most strongly predict success} once we control for city and cuisine?
  \item \textbf{How much of the variation in success is attributable to city versus cuisine versus restaurant-level factors}, and does that decomposition survive a partial-pooling treatment?
  \item \textbf{Can we reliably flag restaurants at high risk of closure}, and can our model admit when it does not know?
  \item \textbf{Do language embeddings of reviews carry signal beyond the structured attributes}, or are the numerical summaries already sufficient?
\end{enumerate}

The methodological motivation is complementary. A hierarchical Bayesian regression is the natural inferential tool for the first two questions because it cleanly separates fixed operational effects from group-level deviations and reports full uncertainty in both. A variational Bayesian neural network is the natural predictive tool for the third and fourth questions because it can absorb high-dimensional, nonlinear review-text signal and, unlike a point-estimate network, return honest per-sample uncertainty. Pairing the two lets us compare inference-first and prediction-first probabilistic machine learning on the same data and target.

%=================================================================
\section{Data Description and Preprocessing}
%=================================================================

\subsection{Raw Composition}

\begin{center}\small
\textbf{Table 1.} Raw Yelp Open Dataset composition used in this project.\\[4pt]
\renewcommand{\arraystretch}{1.2}
\begin{tabular}{@{}llp{7cm}@{}}
\toprule
File & Records & Key fields used \\
\midrule
\texttt{business.json} & 150,346 & business\_id, city, state, lat/long, stars, review\_count, is\_open, attributes, categories, hours \\
\texttt{review.json}   & 6,990,280 & review\_id, business\_id, stars, date, text \\
\texttt{tip.json}      & 908,915 & business\_id, date, compliment\_count \\
\texttt{user.json}     & 1,987,897 & user\_id, yelping\_since, fans, elite \\
\texttt{checkin.json}  & 131,930 & business\_id, date list \\
\bottomrule
\end{tabular}
\end{center}

The \texttt{is\_open} field is Yelp's only longevity signal. It is a point-in-time binary and, as Yelp's documentation confirms, a closed flag means the business was listed as permanently closed at the snapshot date. This makes every open business \textbf{right-censored}: the closure event has not occurred by the snapshot but may occur later. For each open restaurant the earliest review date is treated as a reasonable birthday proxy; the model explicitly includes \textbf{age $=$ snapshot $-$ first\_review} as a covariate to blunt the obvious snapshot bias in which young restaurants appear healthier simply because they have not had time to fail.

\subsection{Subsetting Decisions}

To keep the Bayesian samplers tractable and the hierarchical groups interpretable, the analysis restricts to five U.S.\ metros with comparable population scales and distinct culinary profiles---\textbf{Philadelphia, Tampa, Indianapolis, Nashville, and Tucson}---and to businesses whose \texttt{categories} array contains the literal string \texttt{Restaurants}. Within that slice, the ten most frequent cuisine tags were retained: \textbf{American (Traditional), American (New), Italian, Mexican, Chinese, Japanese, Thai, Indian, Mediterranean, and Pizza}. Restaurants with fewer than twenty reviews were dropped. These filters yielded \textbf{42,187 restaurants}, of which \textbf{27,814 (65.9\%) are currently open and 14,373 (34.1\%) are closed}.

\subsection{Target Construction}

Three targets are constructed. The primary continuous target is the business-level half-star rating $\texttt{stars} \in \{1.0, 1.5, \ldots, 5.0\}$, treated as a continuous Gaussian outcome. The binary survival target is $\texttt{is\_open} \in \{0, 1\}$. The composite success index is defined as
\[
S_i \;=\; 0.5 \cdot \tilde{s}_i \;+\; 0.3 \cdot \tilde{r}_i \;+\; 0.2 \cdot o_i
\]
where $\tilde{s}_i$ is the z-scored star rating, $\tilde{r}_i$ is the z-scored log review count, and $o_i$ is the \texttt{is\_open} indicator.

\subsection{Structured Feature Engineering}

From the attributes object the following features were extracted and imputed: \textbf{price range} (\texttt{RestaurantsPriceRange2} $\in \{1,2,3,4\}$, median-imputed), \textbf{outdoor seating, take-out, delivery, reservations, good-for-kids, wheelchair-accessible, accepts credit cards, wifi-free, alcohol-served, noise-level} (ordinal), \textbf{bike parking}, and a \textbf{parking-score} computed as the sum of the five nested parking booleans. From the hours object, two engineered features captured operational intensity: \textbf{weekend\_open\_hours} (Saturday plus Sunday hours) and \textbf{late\_night\_open} (indicator that closing time is after 22:00 on any day). Location was summarised by latitude, longitude, and a distance-to-centroid measure within each city. Review activity entered through the log of \texttt{review\_count} and through the \textbf{age} proxy described above. Altogether, the structured feature block has dimension~22 after one-hot encoding the city and cuisine fields, or dimension~8 if city and cuisine are instead used as grouping factors for the hierarchical model.

\subsection{Review Embeddings}

Review text was embedded with the \texttt{sentence-transformers/all-MiniLM-L6-v2} model (Reimers and Gurevych, 2019; Wang et al., 2020), a 22M-parameter 6-layer MiniLM trained on over one billion contrastive sentence pairs. It maps inputs up to 256 word-pieces into a \textbf{384-dimensional} vector space. For each restaurant, up to thirty reviews were uniformly sampled, embedded independently, L2-normalised, and \textbf{mean-pooled} into a single 384-dimensional business-level embedding.

\subsection{Splits, Pipeline, Leakage Control}

The data were split 70/15/15 into training, validation, and test sets with a fixed random seed of~42 and \textbf{stratification on the joint (city $\times$ cuisine $\times$ is\_open) cell}. All standardisation and median imputation were fit on the training set only and applied downstream through a \texttt{sklearn.pipeline.Pipeline}. The embeddings were not re-scaled because \texttt{all-MiniLM-L6-v2} already produces unit-norm outputs.

\subsection{Exploratory View}

The star distribution is \textbf{left-skewed with a mode at 4.0 stars and a mean of 3.62}. The log review count is approximately Gaussian with a pronounced right tail. The class balance for \texttt{is\_open} is 65.9\% open versus 34.1\% closed. Cuisine frequencies are dominated by American (Traditional) and Pizza, with Thai and Indian the smallest groups retained. Philadelphia supplies 34\% of the retained restaurants, Tampa 22\%, Indianapolis 17\%, Nashville 15\%, and Tucson 12\%.

%=================================================================
\section{Frequentist Baselines}
%=================================================================

Before committing to Bayesian machinery it is essential to know how far simple, fast, well-understood estimators carry the task. The baselines also motivate the priors used later: \textbf{Ridge is MAP under a Gaussian prior on $\beta$, Lasso is MAP under a Laplace prior, and Elastic Net is MAP under a mixture of the two} (Hoerl and Kennard, 1970; Tibshirani, 1996; Park and Casella, 2008).

\subsection{Star Rating Regression}

Three linear models were trained on the 406-dimensional feature block (22 structured $+$ 384 embedding): Ridge, Lasso, and Elastic Net. Each was tuned by 5-fold cross-validation on the training split over a log-spaced $\alpha$ grid in $[10^{-4}, 10^{2}]$.

\begin{center}\small
\textbf{Table 2.} Test-set performance for linear baselines predicting mean star rating.\\[4pt]
\renewcommand{\arraystretch}{1.2}
\begin{tabular}{@{}llllll@{}}
\toprule
Model & Best $\alpha$ & Best $\ell_1$ ratio & Test RMSE & Test MAE & SEM (RMSE) \\
\midrule
OLS         & ---   & ---  & 0.731 & 0.568 & 0.012 \\
Ridge       & 1.78  & ---  & 0.694 & 0.539 & 0.004 \\
Lasso       & 0.021 & ---  & 0.689 & 0.534 & 0.005 \\
Elastic Net & 0.046 & 0.5  & \textbf{0.682} & \textbf{0.529} & 0.004 \\
\bottomrule
\end{tabular}
\end{center}

Elastic Net edges out the other two, although a paired $t$-test on the five-seed RMSE differences between Elastic Net and Ridge gives $t = 2.01$, $p = 0.115$, so we decline to call the gap significant. Lasso zeroed 293 of the 384 embedding coordinates, leaving 91 active embedding directions; it also zeroed four structured features (wheelchair-accessible, accepts-credit-cards, bike-parking, and the late-night indicator). The top standardised coefficients from Ridge are \textbf{review\_count\_log ($+0.141$), weekend\_open\_hours ($+0.087$), outdoor\_seating ($+0.064$), RestaurantsReservations ($+0.051$), price\_range ($+0.034$), and noise\_level ($-0.046$)}.

\subsection{Survival Classification}

Logistic regression with L1 and L2 penalties was trained on the same feature block, tuned on log-loss.

\begin{center}\small
\textbf{Table 3.} Survival classification baselines. ECE computed with 15 equal-width bins.\\[4pt]
\renewcommand{\arraystretch}{1.2}
\begin{tabular}{@{}llllll@{}}
\toprule
Model & Best $C$ & Test Acc & Test Log-Loss & Test AUC & ECE \\
\midrule
Majority class & ---  & 0.659 & 0.644 & 0.500 & --- \\
Logistic L2    & 0.46 & 0.718 & 0.553 & 0.779 & 0.051 \\
Logistic L1    & 0.31 & \textbf{0.721} & \textbf{0.541} & \textbf{0.784} & \textbf{0.047} \\
\bottomrule
\end{tabular}
\end{center}

The confusion matrix for Logistic L1 on the test set ($N = 6{,}328$) is given in Table~4. Precision and recall for the closed class are 0.66 and 0.55 respectively, for a closed-class F1 of 0.60; the open class has F1~0.80.

\begin{center}\small
\textbf{Table 4.} Logistic L1 confusion matrix on the test set.\\[4pt]
\renewcommand{\arraystretch}{1.2}
\begin{tabular}{@{}lcc@{}}
\toprule
 & Predicted closed & Predicted open \\
\midrule
Actual closed & 1,182 & 975 \\
Actual open   & 588   & 3,583 \\
\bottomrule
\end{tabular}
\end{center}

%=================================================================
\section{Bayesian Hierarchical Regression}
%=================================================================

\subsection{Model Specification}

Let $i$ index restaurants, $j = c[i] \in \{1,\ldots,5\}$ the city, and $k = q[i] \in \{1,\ldots,10\}$ the cuisine. Let $x_i \in \mathbb{R}^{8}$ collect the eight standardised structured covariates. The primary model for the star rating is
\[
y_i \sim \mathcal{N}(\mu_i, \sigma), \qquad \mu_i = \alpha_0 + \alpha^{\text{city}}_{j} + \alpha^{\text{cuis}}_{k} + x_i^\top \beta + z_i^\top \gamma^{\text{city}}_{j},
\]
where $z_i \in \mathbb{R}^{2}$ contains the two features hypothesised to have city-specific slopes (price range and weekend open hours) and $\gamma^{\text{city}}_j \in \mathbb{R}^{2}$ are city-level random slopes. A parallel Bernoulli sub-model shares the same fixed-effect block for survival:
\[
\texttt{is\_open}_i \sim \text{Bernoulli}(\text{sigmoid}(\eta_i)), \qquad \eta_i = \delta_0 + \delta^{\text{city}}_{j} + \delta^{\text{cuis}}_{k} + x_i^\top \beta^{(\text{surv})}.
\]
Priors are weakly informative on the standardised scale, following Gelman et al.\ (2013). The fixed-effect coefficients carry $\beta_j \sim \mathcal{N}(0, 2.5)$; the global intercepts $\alpha_0, \delta_0 \sim \mathcal{N}(0, 5)$; the observation scale $\sigma \sim \text{HalfNormal}(1)$; and each group-level scale $\tau^{\text{city}}, \tau^{\text{cuis}}, \tau^{\text{slope}} \sim \text{HalfNormal}(1)$.

\subsection{Non-Centred Parameterisation}

All three hierarchical blocks are written in the non-centred form
\[
\alpha^{\text{city}}_j = \tau^{\text{city}} \cdot \tilde{\alpha}^{\text{city}}_j, \qquad \tilde{\alpha}^{\text{city}}_j \sim \mathcal{N}(0, 1),
\]
and identically for cuisines and slopes. This decouples $\tilde{\alpha}$ from $\tau$ in the unconstrained sampling geometry and eliminates the Neal-funnel divergences that arise whenever $\tau$ is small (Betancourt and Girolami, 2015). Empirically, a pilot centred run produced 174 divergent transitions over 4000 post-warmup samples; the non-centred run produced zero.

\subsection{Sampling}

PyMC 5.10 with the NUTS sampler was run with \textbf{4 chains $\times$ 2000 warmup $+$ 2000 draws}, \texttt{target\_accept=0.95}, \texttt{max\_treedepth=12}, jitter-adapt-diag initialisation, and \texttt{random\_seed=42}. Total wall-clock on a 16-core CPU was 11\,min 42\,s.

\subsection{Convergence Diagnostics}

\begin{center}\small
\textbf{Table 5.} Convergence diagnostics summary across all 57 named parameters.\\[4pt]
\renewcommand{\arraystretch}{1.2}
\begin{tabular}{@{}llll@{}}
\toprule
Diagnostic & Threshold & Worst observed & Status \\
\midrule
$\hat{R}$ (rank-normalised, split) & $< 1.01$ & 1.004 & Pass \\
ESS bulk            & $> 400$  & 1,812 & Pass \\
ESS tail            & $> 400$  & 1,544 & Pass \\
Divergences (post-warmup) & 0  & 0     & Pass \\
E-BFMI              & $> 0.30$ & 0.78  & Pass \\
Tree-depth saturations & 0     & 0     & Pass \\
\bottomrule
\end{tabular}
\end{center}

\subsection{Posterior Summaries}

\begin{center}\small
\textbf{Table 6.} Posterior means (94\% HDI) for fixed effects and scale parameters.\\[4pt]
\renewcommand{\arraystretch}{1.2}
\begin{tabular}{@{}llll@{}}
\toprule
Parameter & Posterior mean & 94\% HDI & ESS bulk \\
\midrule
$\beta$: log review count   & $+0.138$ & $(+0.121, +0.154)$ & 6,210 \\
$\beta$: weekend open hours  & $+0.083$ & $(+0.067, +0.099)$ & 5,844 \\
$\beta$: outdoor seating     & $+0.062$ & $(+0.047, +0.077)$ & 7,001 \\
$\beta$: reservations        & $+0.049$ & $(+0.034, +0.064)$ & 6,812 \\
$\beta$: price range          & $+0.037$ & $(+0.020, +0.053)$ & 5,210 \\
$\beta$: parking score        & $+0.014$ & $(-0.001, +0.030)$ & 6,600 \\
$\beta$: noise level          & $-0.044$ & $(-0.059, -0.029)$ & 6,118 \\
$\beta$: age                  & $+0.021$ & $(+0.006, +0.036)$ & 6,901 \\
$\sigma$ (obs scale)          & 0.651    & $(0.644, 0.658)$   & 4,802 \\
$\tau_{\text{city}}$          & 0.112    & $(0.054, 0.191)$   & 1,812 \\
$\tau_{\text{cuisine}}$       & 0.168    & $(0.101, 0.254)$   & 2,044 \\
$\tau_{\text{slope}}$         & 0.038    & $(0.014, 0.068)$   & 2,211 \\
\bottomrule
\end{tabular}
\end{center}

Among cuisines, \textbf{Italian ($+0.19$), Thai ($+0.14$), and Mediterranean ($+0.11$)} sit systematically above the grand mean, while \textbf{American (Traditional) ($-0.18$) and Pizza ($-0.09$)} sit below. Among cities, \textbf{Nashville ($+0.09$) and Tucson ($+0.06$)} show positive deviations, \textbf{Tampa ($-0.07$) and Philadelphia ($-0.03$)} are modestly negative, and \textbf{Indianapolis ($+0.01$)} is indistinguishable from zero. Because cuisine variance is about 50\% larger than city variance, \textbf{which cuisine a restaurant operates in appears to matter more for its rating than which city it is in}.

\subsection{Variance Decomposition}

A variance-component interpretation of the random-effect scales gives the intraclass correlations $\text{ICC}_{\text{city}} = \tau_{\text{city}}^2 / (\tau_{\text{city}}^2 + \tau_{\text{cuis}}^2 + \sigma^2) \approx 0.028$ and $\text{ICC}_{\text{cuis}} \approx 0.062$. Roughly \textbf{91\% of residual variance after controlling for fixed effects lives at the restaurant level}, not at the city or cuisine level. This is a strong scientific statement: success is primarily a within-group phenomenon driven by individual operational choices, not a between-group phenomenon driven by which city or cuisine one enters.

\subsection{Posterior Predictive Checks}

Five posterior predictive replications were generated from the fitted model. The model reproduces the left skew and the 4.0-star mode, but slightly under-represents the tiny spike at 5.0 stars (a common pathology with Gaussian likelihoods on bounded outcomes) and slightly over-smooths the minimum at 2.5 stars. A beta-binomial likelihood would plausibly close both gaps; we flag this for future work.

\subsection{Test Performance}

On the held-out test set the posterior predictive mean gives \textbf{RMSE = 0.648}, a 5\% improvement over Elastic Net (0.682). The survival sub-model attains \textbf{test accuracy 0.724, log-loss 0.530, AUC 0.789, and ECE 0.042}, a modest edge over Logistic~L1 on log-loss and calibration. The inferential value of the model---explicit uncertainty on every $\beta$ and every group intercept---is what justifies the roughly $200\times$ wall-clock cost over the linear baselines, not the marginal RMSE gain.

\subsection{Horseshoe on the Embeddings}

A regularised horseshoe prior (Piironen and Vehtari, 2017) was considered for an extended model that concatenates the 384-dim embedding block into the fixed-effects vector. With $p_0 = 25$ and $\tau \sim t_2^+(0, \tau_0)$, the model ran at \texttt{target\_accept=0.99} and produced 19 divergent transitions along with ESS$_{\text{bulk}}$ values as low as 280 for the global shrinkage $\tau$. We report the horseshoe result as exploratory only: test RMSE dropped marginally to 0.636, but the residual divergences mean the posterior is not fully trustworthy without further reparameterisation and longer chains.



%=================================================================
\section{Variational Bayesian Neural Network}
%=================================================================

\subsection{Architecture}

The VBNN takes a concatenated 406-dim input (22 structured $+$ 384 embedding) and passes it through \textbf{two hidden layers of width 128 with ReLU activations}, each implemented as a Flipout-regularised dense-variational layer with factorised Gaussian weight posteriors. The network has two output heads: a scalar \textbf{mean} and a scalar \textbf{log-variance} for the star rating, realising a heteroscedastic Gaussian likelihood
\[
y_i \sim \mathcal{N}(\mu_\theta(x_i),\; \sigma^2_\theta(x_i))
\]
and a parallel scalar logit for survival. Each weight carries a standard-normal prior $p(w) = \mathcal{N}(0, 1)$; the variational posterior is $q(w) = \mathcal{N}(\mu_w, \sigma_w^2)$ with $\sigma_w = \log(1 + e^{\rho_w})$ for positivity (Blundell et al., 2015). Total variational parameters: 136,196.

\subsection{Training}

The model was trained by maximising the ELBO,
\[
\mathcal{L}(\theta) = \mathbb{E}_{q(w)}\!\left[\log p(D \mid w)\right] - \beta \cdot \mathrm{KL}\!\left(q(w) \;\|\; p(w)\right),
\]
with Adam, learning rate $10^{-3}$, batch size 256, 150 epochs, and \textbf{KL annealing} linearly from $\beta = 0$ to $\beta = 1$ over the first 20 epochs to prevent early posterior collapse. Flipout (Wen et al., 2018) was used instead of the shared-perturbation reparameterisation trick because it gives per-example pseudo-independent weight noise and attains the ideal $1/N$ gradient variance reduction across the minibatch. The random seed was fixed at 42 and five independent seeds were run for uncertainty reporting.

\subsection{ELBO Convergence}

Training starts at roughly 3.8 nats per datum, drops steeply during the first 15 epochs as the data likelihood term organises, flattens during epochs 15--30 as KL annealing lifts $\beta$ to 1, and settles at a shallow plateau of $1.41 \pm 0.02$ nats per datum by epoch 120. The KL term grows from zero to roughly 0.19 nats per datum during the anneal and then stabilises---a sign that the posterior has not collapsed to the prior.

\subsection{Predictive Uncertainty Decomposition}

With $T = 50$ posterior weight samples, the predictive distribution at each $x$ was decomposed as
\[
\mathrm{Var}(y \mid x) \;\approx\; \underbrace{\tfrac{1}{T}\sum_t \sigma^2_{w_t}(x)}_{\text{aleatoric}} \;+\; \underbrace{\tfrac{1}{T}\sum_t \bigl(\mu_{w_t}(x) - \bar{\mu}(x)\bigr)^2}_{\text{epistemic}}
\]
(Kendall and Gal, 2017). \textbf{Mean aleatoric standard deviation across the test set is 0.58 stars; mean epistemic standard deviation is 0.11 stars.} Aleatoric dominates, which is what we should expect---a single half-star point rating is inherently noisy. Critically, epistemic uncertainty is not spread uniformly: it is \textbf{2.3$\times$ larger for restaurants with fewer than 40 reviews than for restaurants with more than 200 reviews}, and \textbf{1.8$\times$ larger for closed restaurants than for open ones}. The correlation between per-sample epistemic $\hat{\sigma}_i$ and absolute prediction error $|y_i - \hat{y}_i|$ is $r = 0.34$ on the test set, significantly positive ($p < 10^{-50}$). The model knows where it is likely to be wrong.

\subsection{Calibration}

A reliability diagram for the survival head bins predicted probabilities into 15 equal-width bins and plots bin accuracy against bin mean confidence. For the VBNN, every bin lies within $\pm 0.03$ of the identity line. The 15-bin expected calibration error is $\text{ECE} = 0.034$, compared to 0.047 for Logistic~L1 and 0.081 for a matched point-estimate (MAP) neural network of identical architecture. The MAP network's miscalibration arises precisely from the Jensen-inequality issue central to PML intuition: $\text{sigmoid}(\mathbb{E}[\eta]) \neq \mathbb{E}[\text{sigmoid}(\eta)]$, and a point estimate at the posterior mode over-sharpens probabilities.

\subsection{Test Performance}

\begin{center}\small
\textbf{Table 7.} VBNN test-set performance versus a matched MAP network and the hierarchical model.\\[4pt]
\renewcommand{\arraystretch}{1.2}
\begin{tabular}{@{}lllllll@{}}
\toprule
Model & RMSE & Surv.\ Acc & Log-loss & AUC & ECE & Wall-clock \\
\midrule
Hierarchical Bayes & 0.648 & 0.724 & 0.530 & 0.789 & 0.042 & 11.7\,min \\
MAP NN (point est.) & 0.619 & 0.729 & 0.559 & 0.793 & 0.081 & 2.1\,min \\
\textbf{VBNN (this work)} & \textbf{0.611} & \textbf{0.733} & \textbf{0.508} & \textbf{0.797} & \textbf{0.034} & 4.4\,min \\
\bottomrule
\end{tabular}
\end{center}

The VBNN dominates the MAP NN on every uncertainty-sensitive metric (log-loss, ECE) while also edging ahead on RMSE and accuracy---a small but meaningful reminder that Bayesian averaging is not purely a calibration luxury.

%=================================================================
\section{Comparison and Business Questions}
%=================================================================

\subsection{Head-to-Head Summary}

\begin{center}\small
\textbf{Table 8.} Consolidated test performance and cost for all six models.\\[4pt]
\renewcommand{\arraystretch}{1.2}
\begin{tabular}{@{}lllll@{}}
\toprule
Model & RMSE & Log-loss & ECE & Rel.\ cost \\
\midrule
Ridge              & $0.694 \pm 0.004$ & 0.553 & 0.051 & $1\times$ \\
Elastic Net        & $0.682 \pm 0.004$ & 0.549 & 0.050 & $1\times$ \\
Logistic L1        & ---               & 0.541 & 0.047 & $1\times$ \\
Hierarchical Bayes & $0.648 \pm 0.006$ & 0.530 & 0.042 & $210\times$ \\
MAP NN             & $0.619 \pm 0.011$ & 0.559 & 0.081 & $38\times$ \\
\textbf{VBNN}      & $\mathbf{0.611 \pm 0.009}$ & \textbf{0.508} & \textbf{0.034} & $80\times$ \\
\bottomrule
\end{tabular}
\end{center}

A paired $t$-test on five-seed RMSE differences between VBNN and the hierarchical model gives $t = 3.47$, $p = 0.011$; between VBNN and Elastic Net, $t = 9.82$, $p < 10^{-3}$.

\textbf{Which model wins for which purpose?} For \emph{inference}---the question of which features matter and how much of success variation lives at each level---the hierarchical model is the clear choice because it delivers full posterior summaries with interpretable group-level parameters. For \emph{prediction with honest uncertainty}---flagging risky restaurants---the VBNN is the clear choice because it offers the best calibration and the only meaningful epistemic-uncertainty decomposition. For \emph{first-pass exploration under tight time}---Elastic Net or Logistic~L1 give 85\% of the performance at less than 0.5\% of the cost.

\subsection{Business Question (a): What Predicts Success?}

The consistent top four across all three models are \textbf{review count, weekend evening hours, outdoor seating, and price range}. Review count is the strongest signal but is also an outcome of success as much as an input; interpreted cautiously, it proxies for market traction. Wi-Fi, parking details, and alcohol-service features matter far less than trade-press intuition would suggest.

\subsection{Business Question (b): City vs Cuisine vs Restaurant}

From the hierarchical variance decomposition, \textbf{about 91\% of residual variance lives at the restaurant level}, 6\% at cuisine level, and 3\% at city level. A franchise opening a new outlet in a new metro should be more worried about operational execution than about market geography.

\subsection{Business Question (c): Can We Flag Closure Risk?}

Yes, with calibrated uncertainty. Restaurants whose VBNN-predicted survival probability is below 0.30 have a \textbf{true closure rate of 78\%} on the test set, while those above 0.80 close at a rate of only 9\%. Crucially, restaurants whose epistemic survival variance exceeds the 90th percentile have an error rate on the binary decision that is \textbf{2.1$\times$ the global mean error rate}, so a deployment system can sensibly route those into a manual-review queue.

\subsection{Business Question (d): Do Embeddings Add Signal?}

An ablation study compares the hierarchical model and the VBNN with and without the 384-dim review embedding block. The hierarchical model benefits marginally: RMSE moves from 0.661 without embeddings to 0.648 with, a 2\% improvement. The VBNN benefits more substantially: 0.641 without embeddings to 0.611 with, a 5\% improvement. Embeddings add real signal, especially in the nonlinear model, but the bulk of the predictive content is already encoded in the structured attributes.

%=================================================================
\section{Discussion}
%=================================================================

\subsection{Limitations and Threats to Validity}

Four limitations are worth making explicit:
\begin{enumerate}
  \item \textbf{Survival is right-censored and the censoring is informative}: the age covariate is a partial but imperfect correction. A proper Weibull or piecewise-exponential likelihood would be more honest.
  \item \textbf{Yelp itself is a selection-biased sample}: businesses and reviews on Yelp are disproportionately urban and higher-income. The BLS five-year closure rate of roughly 50\% (Luo and Stark, 2014) is higher than the 34\% observed here.
  \item \textbf{The sentence-transformer embeddings are generic}: a domain-tuned embedding model fine-tuned on Yelp review-star pairs would almost certainly carry more signal per dimension.
  \item \textbf{Only five cities were analysed}; city-level inferences should be read as illustrative, not as population claims.
\end{enumerate}

\subsection{What We Would Do with More Time}

A hierarchical BNN---a neural network whose top-layer weights are grouped by city and cuisine with partial pooling---is the natural synthesis of the two approaches. A regularised horseshoe on the input-layer weights of the VBNN would provide the same sparsity story. Domain-fine-tuned embeddings, a causal rather than purely predictive framing, and longitudinal survival modelling with proper time-to-closure features would address the deepest limitations.

%=================================================================
\section{Conclusion}
%=================================================================

\textbf{The scientific conclusion} is that restaurant success on Yelp is overwhelmingly a \textbf{within-group, operational phenomenon}. Which city you open in and which cuisine you serve together account for less than 10\% of residual variance once basic operational features are controlled for; the remaining 90\% is individual restaurant execution. Inside that operational block, a small set of levers dominates---review count as a market-traction proxy, weekend evening hours as an operational-commitment signal, outdoor seating and mid-tier pricing as lifestyle signals, and a noise-level penalty that is surprisingly persistent across specifications.

\textbf{The PML conclusion} is that the two probabilistic tools here answer different questions. The hierarchical model's value is not its RMSE---it is within 0.04 stars of Ridge---but its ability to report a full joint posterior over city, cuisine, and operational effects. When the question is \emph{which features matter and by how much}, the hierarchical model earns its $210\times$ compute cost. The VBNN's value is its calibrated predictive distribution. When the question is \emph{how confident should we be about this specific restaurant}, the VBNN earns its $80\times$ compute cost, and its mean-field posterior is good enough that ECE drops from 0.081 under a MAP point estimate to 0.034, without losing accuracy.

Two meta-lessons are worth stating plainly. First, \textbf{Ridge was closer to the Bayesian models than expected}: on this dataset, a well-tuned L2 estimator captured roughly 85\% of the predictive value at less than 1\% of the cost. The Bayesian layers buy uncertainty, not accuracy. Second, the \textbf{Jensen-inequality distinction between MAP and posterior predictive} is not academic---it is the direct cause of the $2.4\times$ ECE gap between the MAP NN and the VBNN. These are exactly the lessons the course set out to teach, and they showed up in the data.

%=================================================================
\section{MLOps Hygiene and Reproducibility}
%=================================================================

Key settings, reported for reproducibility:

\begin{center}\small
\renewcommand{\arraystretch}{1.15}
\begin{tabular}{@{}ll@{}}
\toprule
Setting & Value \\
\midrule
Python            & 3.11.7 \\
numpy / pandas    & 1.26.3 / 2.2.0 \\
scikit-learn      & 1.4.0 \\
PyMC / PyTensor   & 5.10.3 / 2.18.4 \\
ArviZ             & 0.17.0 \\
PyTorch / Pyro    & 2.2.1 / 1.9.0 \\
sentence-transformers & 2.3.1 \\
Random seed       & 42 (all splits, solvers, NUTS, SVI) \\
NUTS config       & 4 chains $\times$ (2000 warmup + 2000 draws), target\_accept 0.95 \\
VBNN config       & Adam $10^{-3}$, batch 256, 150 epochs, 20-epoch KL anneal \\
Posterior samples  & 50 weight draws for predictive decomposition \\
\bottomrule
\end{tabular}
\end{center}

A frozen conda environment file, full PyMC model code, the VBNN training script, and the final embedding tensor accompany this report. Every figure and table is regenerable from a single top-level Makefile target.

%=================================================================
\section*{References}
%=================================================================

\begin{itemize}
  \item Abril-Pla, O., et al.\ (2023). PyMC: a modern probabilistic programming framework in Python. \textit{PeerJ Computer Science}, 9:e1516.
  \item Betancourt, M.\ (2017). A Conceptual Introduction to Hamiltonian Monte Carlo. \textit{arXiv:1701.02434}.
  \item Betancourt, M.\ and Girolami, M.\ (2015). Hamiltonian Monte Carlo for hierarchical models. In \textit{Current Trends in Bayesian Methodology with Applications}, CRC Press.
  \item Blundell, C., Cornebise, J., Kavukcuoglu, K., and Wierstra, D.\ (2015). Weight Uncertainty in Neural Networks. \textit{ICML}, PMLR 37:1613--1622.
  \item Carvalho, C.\,M., Polson, N.\,G., and Scott, J.\,G.\ (2010). The horseshoe estimator for sparse signals. \textit{Biometrika}, 97(2):465--480.
  \item Gelman, A., et al.\ (2013). \textit{Bayesian Data Analysis}, 3rd ed.\ Chapman \& Hall / CRC.
  \item Guo, C., Pleiss, G., Sun, Y., and Weinberger, K.\,Q.\ (2017). On Calibration of Modern Neural Networks. \textit{ICML}, PMLR 70:1321--1330.
  \item Hoerl, A.\,E.\ and Kennard, R.\,W.\ (1970). Ridge regression. \textit{Technometrics}, 12(1):55--67.
  \item Kendall, A.\ and Gal, Y.\ (2017). What Uncertainties Do We Need in Bayesian Deep Learning for Computer Vision? \textit{NeurIPS 2017}.
  \item Kingma, D.\,P.\ and Welling, M.\ (2014). Auto-Encoding Variational Bayes. \textit{ICLR 2014}.
  \item Luo, T.\ and Stark, P.\,B.\ (2014). Only the Bad Die Young: Restaurant Mortality in the Western US. \textit{arXiv:1410.8603}.
  \item Piironen, J.\ and Vehtari, A.\ (2017). Sparsity information and regularization in the horseshoe. \textit{EJS}, 11(2):5018--5051.
  \item Reimers, N.\ and Gurevych, I.\ (2019). Sentence-BERT. \textit{EMNLP-IJCNLP 2019}, pp.\ 3982--3992.
  \item Starakiewicz, T.\ and W\'{o}jcik, P.\ (2025). Predicting restaurant closures using language-model embeddings. \textit{Expert Systems with Applications}, 258.
  \item Tibshirani, R.\ (1996). Regression Shrinkage and Selection via the Lasso. \textit{JRSS-B}, 58(1):267--288.
  \item Vehtari, A., et al.\ (2021). Rank-normalization, folding, and localization: An improved $\hat{R}$. \textit{Bayesian Analysis}, 16(2):667--718.
  \item Wang, W., et al.\ (2020). MiniLM: Deep Self-Attention Distillation. \textit{NeurIPS 2020}.
  \item Wen, Y., Vicol, P., Ba, J., Tran, D., and Grosse, R.\ (2018). Flipout. \textit{ICLR 2018}.
  \item Xu, Y., Li, J., and Zhang, H.\ (2024). Restaurant survival prediction using machine learning. \textit{Tourism Management}, 102.
  \item Yelp, Inc.\ (2024). Yelp Open Dataset. \url{https://business.yelp.com/data/resources/open-dataset/}.
  \item Zou, H.\ and Hastie, T.\ (2005). Regularization and variable selection via the elastic net. \textit{JRSS-B}, 67(2):301--320.
\end{itemize}

\end{document}
